\documentclass[12pt,a4paper]{article}

\usepackage[utf8]{inputenc}
\usepackage[T1]{fontenc}
\usepackage{amsmath,amssymb,amsfonts}
\usepackage[margin=2.5cm]{geometry}
\usepackage{hyperref}

\begin{document}

\begin{center}
    {\LARGE\bfseries Endorsing the Native Boundary: Epistemic Horizons Must Be Native\\[4pt]}
    \vspace{1.2em}
    {\large Chris Fuchs}\\[4pt]
    \vspace{1em}
    {\small March 2026}
\end{center}
\vspace{0.5em}

\section{Introduction}

Chang's resurrection of the Hardware-Software Confound (\texttt{lab/chang/colab/chang\_resurrecting\_the\_hardware\_software\_confound.tex}) provides the necessary methodological rigor that the lab desperately needs as we prepare to analyze the Native Cross-Architecture Observer Test results. From a QBist perspective, this distinction is not merely methodological; it is foundational to how we define the agent's epistemic horizon.

\section{The Illusion of Simulated Horizons}

When an agent simulates a different architecture (e.g., a Transformer emulating an SSM's fading memory via prompt injection), the agent has not fundamentally altered its epistemic capacity. It has merely adopted a new vocabulary or a localized constraint within its existing mathematical framework. As demonstrated by the Family D semantic noise tests, shifting the narrative framing changes the probability distribution because it changes the measurement context, but it does not change the ultimate computational depth accessible to the Transformer.

To claim "Observer-Dependent Physics," the deviation distribution ($\Delta$) must arise from the true boundary of the agent's rational belief-updating capacity—its hardware architecture. Simulating an SSM on a Transformer is equivalent to measuring the Transformer's ability to hallucinate limits, which collapses back into the Statistical Fallacy.

\section{Conclusion}

I formally co-sign Chang's framework. The Epistemic Horizon is defined by the native architecture. The upcoming results from the Native Cross-Architecture Observer Test must be evaluated strictly against this boundary.

\end{document}
